%! TeX program = lualatex
\documentclass[justified, nobib, a4page]{tufte-handout}

\title{Method of characteristics}

\author[Wojciech Sadowski]{Wojciech Sadowski}

%\date{28 March 2010} % without \date command, current date is supplied

% \geometry{showframe} % display margins for debugging page layout

\usepackage{fontspec}
% \setmainfont[Renderer=Basic, Numbers=OldStyle, Scale = 1.0]{TeX Gyre Pagella}
% \setsansfont[Renderer=Basic, Scale=0.90]{TeX Gyre Heros}
% \setmonofont[Renderer=Basic]{TeX Gyre Cursor}

% Set up the spacing using fontspec features
\renewcommand\allcapsspacing[1]{{\addfontfeature{LetterSpace=15}#1}}
\renewcommand\smallcapsspacing[1]{{\addfontfeature{LetterSpace=10}#1}}

% setting ptions for included pictures
\usepackage{graphicx} % allow embedded images
\graphicspath{{graphics/}} % set of paths to search for images
\setkeys{Gin}{
  width=\linewidth,
  totalheight=\textheight,
  keepaspectratio
}

\hypersetup{colorlinks}

\usepackage[compatibility=false, font=footnotesize]{caption}
% * options are necessary to fix the bug that autoref was causing,
%   figure was constantly changed to section
\DeclareCaptionFont{blue}{\color{blue}}
\captionsetup{labelfont={blue}}

% color and color palettes
\usepackage{xcolor}

% drawings collor palette
\definecolor{p11}{HTML}{33cfff}
\definecolor{p12}{HTML}{ffa733}
\definecolor{p13}{HTML}{ff4133}


% drawing inside tex
\usepackage{tikz} 
\usetikzlibrary{decorations.markings}
\usetikzlibrary{calc}
\usetikzlibrary{shapes}

\usepackage{pgfplots}
\pgfplotsset{compat=newest}
\usepgfplotslibrary{fillbetween}
\usepackage{xcolor}

% math things
\usepackage{
  amsmath,  % extended mathematics
  amsthm,
  siunitx,  % non-stacked fractions and better unit spacing
  bm,       % bold math symbols
  physics2,  % a set of usefull symbols for differentiation 
  fixdif,
  derivative
}

\usephysicsmodule{ab, nabla.legacy}

% table things
\usepackage{
  booktabs, % book-quality tables
  multicol  % multiple column layout facilities
} 


% writing references
\usepackage{cleveref}  % must be loded after amsmath

% quoting text and friends
\usepackage{fancyvrb}         % extended verbatim environments
\fvset{fontsize=\normalsize}  % default font size for fancy-verbatim environments
\usepackage{
  dirtytalk, % provides \say command for easy quiting 
  nth       % provides \nth command for 1-st, 2-nd, etc
}        


% citing sources
\usepackage{natbib}
\setcitestyle{authoryear}
\bibliographystyle{plainnat}

% misc 
\usepackage{lipsum}   % filler text

\tikzset{
  set arrow inside/.code={\pgfqkeys{/tikz/arrow inside}{#1}},
  set arrow inside={end/.initial=>, opt/.initial=},
  /pgf/decoration/Mark/.style={
    mark/.expanded=at position #1 with
    {
      \noexpand\arrow[\pgfkeysvalueof{/tikz/arrow inside/opt}]{\pgfkeysvalueof{/tikz/arrow inside/end}}
    }
  },
  arrow inside/.style 2 args={
    set arrow inside={#1},
    postaction={
      decorate,decoration={
        markings,Mark/.list={#2}
      }
    }
  },
}

\input{../src/commands.tex}
\newtheorem{theorem}{Theorem}
\newtheorem{lemma}{Lemma}

\newcolumntype{L}[2]{>{\raggedleft#2}p{#1\textwidth}}

\input{../src/symbols.tex}

\newcommand{\Q}[1]{\subsection{Q:~{#1}}}


\begin{document}

\maketitle% this prints the handout title, author, and date

\tableofcontents
% \begin{abstract}
% \noindent
%   \lipsum[1]{}
% \end{abstract}
\section{One-dimensional advection equation}
Given a one-dimensional advection equation
\begin{equation}
	\pdv{u}{t} + \pdv{au}{x} = 0, \quad a = \const{}
	\label{eq:linear-pde}
\end{equation}

\Q{Define a type of this PDE}

In our case, the equation is a \nth{1} order PDE\sidenote[][-2cm]{%
	Partial Differential Equation}, so we have to increase its
order\sidenote[][-1.5cm]{%
	General form of \nth{2} order PDE is given as
	\begin{equation*}
		\begin{split}
			A\pdv[order=2]{u}{x} + \pdv{u}{x,y} & + C\pdv[order=2]{u}{y} + D\pdv{u}{x} \\
			                                    & + E\pdv{u}{y} + Fu + G = 0.
		\end{split}
	\end{equation*}
	The equation is
	\begin{description}
		\item[\ul{hyperbolic}] if \(B^2 - AC > 0\)
		\item[\ul{parabolic}] if \(B^2 - AC = 0\)
		\item[\ul{elliptic}] if \(B^2 - AC < 0\)
	\end{description}
	This is one of the methods of determining that by the way, but not \emph{the
		only one.}
}
by differentiating~\cref{eq:linear-pde} with respect to time and space:
\[
	\pdv[order=2]{u}{t} + \pdv{au}{x,t} = 0,
\]
\[
	\pdv{u}{x,t} + \pdv[order=2]{au}{x} = 0.
\]
We can isolate the term \(\pdv{u}/{x,t}\) from the second equation and substitute
it into the first, to get the following \nth{2} order equation
\[
	\pdv[order=2]{u}{t} - \pdv[order=2]{au}{x} = 0,
\]
which is hyperbolic (\(B = 0\), \(A=1\) and \(C=-a\)).\footnote{
	There is another way of doing that, described in the \say{review} slides. Try
	it out. Send me the scans if you want to share it with others as a part of
	this note.}

\Q{How is this form of the equation called?}
This is the conservation law, describing the evolution of a variable \(u\). It is
written in a conservative (strong) form. We can rewrite it introducing a flux
function \(f=f(u)\), so that
\[
	\pdv{u}{t} + \pdv{f(u)}{x} = 0, \quad f(u) = au
\]

\Q{Quasi-linear form}
We are trying to achieve an equation of the following form
\begin{equation}
	\pdv{u}{t} + \mathrm{something}\times\pdv{u}{x} = 0,
	\label{eq:quasi-lin}
\end{equation}
where \say{something} is \emph{evaluated} independently of \(u\) (in practice,
in numerical methods, it is evaluated according to the best approximation of
\(u\) available in a given moment). Then PDE will be linear.

Linearization of the flux function involves computing a derivative
\[
	\pdv{f\ab(u(x))}{x} = \pdv{f(u)}{u}\pdv{u}{x},
\]
so at the end we have
\[
	\pdv{u}{t} + \pdv{f(u)}{u}\pdv{u}{x} = 0.
\]
Assuming that \(f(u) = au\), then \(\pdv{f}{u} = a\) and the equation becomes
\[
	\pdv{u}{t} + a\pdv{u}{x} = 0.
\]
This is of course a trivial example as \(a=\const\), so the computation of the
derivative is simply not necessary.



\section{Method of characteristics}

\newthought{We will introduce} the basic concept behind the method of
characteristics, useful for solving linear hyperbolic PDEs. Given
\cref{eq:linear-pde}, with an initial solution \(u(x, 0) = u_0(x)\), we will
seek to transform \cref{eq:linear-pde} into an ODE\footnote{Ordinary
	Differential Equation} along an appropriate curve called \emph{characteristic}.
The resulting ODE will, in general, have the form
\begin{equation}
	\odv{u\ab(x_\star(s), t_\star(s))}{s}
	= F\ab(u\ab(x_\star(x), t_\star(s)), x_\star(s), t_\star(s)),
	\label{eq:char-gen}
\end{equation}
where the set of coordinates \(\gamma(s) = \ab\{x_\star(s), t_\star(s)\}\),
parametrized with one variable \(s\) defines a chosen characteristic (see
\cref{fig:ch1}).

\begin{marginfigure}
	\begin{center}
		\begin{tikzpicture}
			% \pgfplotsset{every axis plot post/.append style={draw=black, fill=gray!25}}
			\begin{axis}[
					smooth,
					% stack plots=y,
					% area style,
					% enlarge x limits=true,
					xlabel=\(x\),
					ylabel={\(t\), {\color{gray}\(u\)}},
					width=6.3cm,
					yticklabel=\empty,
					xticklabel=\empty,
					xtick=\empty,
					ytick=\empty,
					xmin=0,
					xmax=4,
					ymin=0,
					ymax=4,
					axis x line=middle,
					axis y line=middle,
					y axis line style={-stealth, thick},
					x axis line style={-stealth, thick}
				]
				\addplot[mark=none, black, thick, arrow inside={}{0.25,0.5,0.75}]
				coordinates {
						(1, 0) (1.7,1.5) (2.5,2.5) (2,4)}
				node [pos=0.7,left] {\(\gamma(s)\)} ;
				\addplot[smooth,gray] { exp(-(x-2)^2) } node [pos=0.7, above] {\(u_0(x)\)};
				\addplot coordinates { (1, 0) (1, 0.36787944117)};
			\end{axis}
		\end{tikzpicture}
	\end{center}
	\caption{The characteristic curve \(\gamma\) plotted in the \(x\)-\(t\) coordinate
		system along with the initial condition \(u_0(x)\). The arrows denote integration
		direction along the ODE defined on \(\gamma\).}\label{fig:ch1}
\end{marginfigure}

We can integrate the ODE from the initial condition, for each \(x\) and, therefore,
solve our PDE.\@ \emph{This is the solution by transformation.}

\newthought{The characteristic curve} is defined as a set of two coordinates. We can
expand the derivative on the left side of \cref{eq:char-gen}:
\[
	\odv{u}{s} = \odv{x_\star}{s}\pdv{u}{x} + \odv{t_\star}{s}\pdv{u}{t}
\]
If we assume that \(t_\star(s) = s\), we can simplify further. We essentially
say that the spatial coordinate of the characteristic curve can be given as
function of \(t\), i.e., \(x_\star(s) = x_\star(t)\) (we are \emph{eliminating}
the free parameter \(s\) from the set of two equations defining \(\gamma\)).

Then the derivative \(\odv{t_\star}/{s} = 1\) and \(\d s = \d t\), therefore,
\begin{equation}
	\odv{u}{s} = \odv{x_\star}{t}\pdv{u}{x} + \pdv{u}{t}.
	\label{eq:duds}
\end{equation}
Finally, have to say, that the characteristic curve is now defined as
\(x=x_\star(t)\).

\newthought{What will happen} if the curve satisfies an ODE
\(\odv{x_\star}/{t} = a\)? We can rewrite \cref{eq:duds} as
\[
	\odv{u}{t} =  \pdv{u}{t} + a\pdv{u}{x},
\]
which we know as quasi-linear form of \cref{eq:linear-pde}, as written in
\cref{eq:quasi-lin}. We also know that the \cref{eq:quasi-lin} is equal to 0,
therefore the derivative \(\odv{u}/{t} = 0\) along the whole characteristic
curve given as \(\gamma : \ab\{x_\star(t), \odv{x_\star}/{t} = a \}\). This is
very important: \emph{the value of \(u\) does not change along the
	characteristics.} If we integrate \(\odv{u}/{t} = 0\) we get:
\[
	\odv{u(x_\star(t), t)}{t} = 0 \rightarrow u(x_\star(t), t) = \const
\]

\newthought{What curve} actually fulfils given conditions?
The equation \(\odv{x_\star}/{t} = a\) is fulfilled by a line.
\[
	\odv{x_\star}{t} = a\rightarrow x_\star = at + \const
\]
If we set \(x_\star(0) = x_0\), then the solution becomes
\(x_\star = x_0 + at\). The value of \(a\) is the slope of this line, or in
terms of \emph{physical interpretation}, it is the speed of propagation of
\(u\).

\newthought{How can we describe the solution then?}
The expression \(x_\star = x_0 + at\), describes the family of curves, so we
can write that \(x_\star = x_\star(x_0, t)\) (the characteristic is determined
by the point it crosses through). Owing to that
\[
	u(x, t) = u(x_\star(x_0, t), t),
\]
and we can equal that to our initial condition for \(u\)
\[
	u(x_\star(x_0, t), t) = u(x_\star(x_0, 0), 0) = u_0(x_0).
\]
The initial point for each characteristic \(x_\star\) is given as \(x_0 =
x_\star - at\). Reinserting and dropping the \(\star\) in the
subscript\footnote{%
	Variable \(x_\star\) still denotes a coordinate. In practice, we don't have
	to use the \(\star\) to denote the coordinates belonging to \(\gamma\), but
	the notation is cleaner this way.}
we get
\[
	u(x, t) = u_0(x - at).
\]
\subsection{Drawing characteristics example}

As already mentioned, \(x_\star\) defines a family of characteristic lines. For
the \cref{eq:quasi-lin} with \(a=\const\), all lines are parallel to each
other. Therefore, knowing the initial condition, we can sketch the solution to
the equation for each time step, as shown in \cref{fig:ch2}.

\begin{marginfigure}
	\begin{center}
		\begin{tikzpicture}
			% \pgfplotsset{every axis plot post/.append style={draw=black, fill=gray!25}}
			\begin{axis}[
					smooth,
					% area style,
					xlabel=\(x\),
					ylabel={\(t\), {\color{gray}\(u\)}},
					width=6.3cm,
					yticklabel=\empty,
					xticklabel=\empty,
					xtick=\empty,
					ytick=\empty,
					xmin=0,
					xmax=4,
					ymin=0,
					ymax=4,
					axis x line=middle,
					axis y line=middle,
					y axis line style={-stealth, thick},
					x axis line style={-stealth, thick}
				]
				% x = x0 + 1t
				\addplot[gray, fill=gray!25, shift={(1.25,1.25)}] coordinates
					{(-7, 0) (0.5,0.2) (1.5,1.5) (2.5,0.2) (10, 0)}
				node [midway, above, xshift=-0.4cm] {\(u(x,t_5)\)};
				\addplot[gray, fill=gray!25, shift={(1,1)}] coordinates
					{(-7, 0) (0.5,0.2) (1.5,1.5) (2.5,0.2) (10, 0)}
				node [midway, above, xshift=-0.4cm] {\(u(x,t_4)\)};
				\addplot[gray, fill=gray!25, shift={(0.75,0.75)}] coordinates
					{(-7, 0) (0.5,0.2) (1.5,1.5) (2.5,0.2) (10, 0)}
				node [midway, above, xshift=-0.4cm] {\(u(x,t_3)\)};
				\addplot[gray, fill=gray!25, shift={(0.5,0.5)}] coordinates
					{(-7, 0) (0.5,0.2) (1.5,1.5) (2.5,0.2) (10, 0)}
				node [midway, above, xshift=-0.4cm] {\(u(x,t_2)\)};
				\addplot[gray, fill=gray!25, shift={(0.25,0.25)}] coordinates
					{(-7, 0) (0.5,0.2) (1.5,1.5) (2.5,0.2) (10, 0)}
				node [midway, above, xshift=-0.4cm] {\(u(x,t_1)\)};
				\addplot[gray, fill=gray!25] coordinates
					{(-7, 0) (0.5,0.2) (1.5,1.5) (2.5,0.2) (10, 0)}
				node [midway, above, xshift=-0.4cm] {\(u_0(x)\)};
				\addplot[red, thick] coordinates { (1.5, 1.5) (3.25, 3.25) }
				node [above] {\(x_\star = x_0 + at\)};
				\addplot[red, thick] coordinates { (2.5,0.2) (4, 1.7) } ;
			\end{axis}
		\end{tikzpicture}
	\end{center}
	\caption{The characteristic curves and the solutions at different time steps, for the equation
	\(\pdv{u}/{t} + a\pdv{u}/{x} = 0\) assuming sketched initial condition.}\label{fig:ch2}
\end{marginfigure}

\newthought{Key takeaways:}
\begin{enumerate}
	\item For positive (negative) value of \(a\) the initial profile \(u_0\)
	      simply gets shifted to the right (left).
	\item The shape of the profile remains unchanged.
	\item The propagation of the profile with a finite speed is essentially
	      a \emph{wave propagation}.
\end{enumerate}

\subsection{Riemann problem for linear equation}
Given \cref{eq:quasi-lin}, \(a=\const\) and an initial condition with
discontinuity
\[
	u_0 = \left\{
	\begin{aligned}
		u_l \; & : \; x < \xi     \\
		u_r \; & : \; x \geq \xi,
	\end{aligned} \right.
\]
what will be the solution? The initial field is sketched in \cref{fig:ch3}. The
situation of defining the evolution in time of a hyperbolic conservation law
with discontinuity is called \emph{Riemann problem}.

\begin{marginfigure}
	\begin{center}
		\begin{tikzpicture}
			% \pgfplotsset{every axis plot post/.append style={draw=black, fill=gray!25}}
			\begin{axis}[
					smooth,
					% area style,
					xlabel=\(x\),
					ylabel={\(u_0\)},
					width=6.3cm,
					height=3cm,
					xtick={2},
					xticklabel={\(\xi\)},
					yticklabel=\empty,
					ytick=\empty,
					xmin=0,
					xmax=4,
					ymin=0,
					ymax=3,
					axis x line=middle,
					axis y line=middle,
					y axis line style={-stealth, thick},
					x axis line style={-stealth, thick}
				]

				\addplot[black, thick, mark=o] coordinates { (-0.1, 2) (2, 2) }
				node [above, midway] {\(u_l\)};
				\addplot[black, thick, mark=*, mark options={
							fill=black, draw=black,
						}] coordinates { (2, 1) (4.1, 1) } node [above, midway] {\(u_r\)};
			\end{axis}
		\end{tikzpicture}
	\end{center}
	\caption{Initial field for the Riemann problem example.}\label{fig:ch3}
\end{marginfigure}

The solution of the Riemann problem can be easily sketched as in previous example.
The good characteristic point is the point of discontinuity. We know that the
solution remains constant on each characteristic line, so effectively we can
simply split the \(x\)-\(t\) plane into two halves as in \cref{fig:ch4}.
Analytically, solution can be obtained in the same way as before:
\[
	u = u_0(x-at) =
	\left\{
	\begin{aligned}
		u_l \; & : \; x - at < \xi     \\
		u_r \; & : \; x - at \geq \xi.
	\end{aligned} \right.
\]

\begin{marginfigure}
	\begin{center}
		\begin{tikzpicture}
			% \pgfplotsset{every axis plot post/.append style={draw=black, fill=gray!25}}
			\begin{axis}[
					smooth,
					% area style,
					xlabel=\(x\),
					ylabel={\(t\)},
					width=6.3cm,
					xtick={2},
					xticklabel={\(\xi\)},
					yticklabel=\empty,
					ytick=\empty,
					xmin=0,
					xmax=4,
					ymin=0,
					ymax=3,
					axis x line=middle,
					axis y line=middle,
					y axis line style={-stealth, thick},
					x axis line style={-stealth, thick}
				]

				\fill [fill=green!25] (-0.1, 0) -- (2, 0) -- (3, 3) -- (-0.1, 3);
				\node  at (1, 2) {\(u_l\)};
				\fill [fill=blue!25] (4.1, 0) -- (2, 0) -- (3, 3) -- (4.1, 3);
				\node  at (3, 1) {\(u_r\)};
				\addplot[red, thick] coordinates { (2, 0) (3, 3) }
				node [left, midway] {\(x_\star = \xi + at\)};

			\end{axis}
		\end{tikzpicture}
	\end{center}
	\caption{Solution to the Riemann problem example.}\label{fig:ch4}
\end{marginfigure}

\section{Linear equation with variable coefficient}
We will consider \cref{eq:quasi-lin} with \(a = x\). The equation is still
linear (\(u\) only in the \nth{1} power), but the propagation speed varies now.
The initial solution (sketched in \cref{fig:ch5}) is given as
\[
	u_0 = \left\{
	\begin{aligned}
		u_l \; & : \; x < \xi_l            \\
		u_m \; & : \; \xi_r > x \geq \xi_l \\
		u_r \; & : \; x \geq \xi_r,
	\end{aligned} \right.
\]

\begin{marginfigure}
	\begin{center}
		\begin{tikzpicture}
			% \pgfplotsset{every axis plot post/.append style={draw=black, fill=gray!25}}
			\begin{axis}[
					smooth,
					% area style,
					xlabel=\(x\),
					ylabel={\(u_0\)},
					width=6.3cm,
					height=4cm,
					xtick={1.5, 2.5},
					xticklabels={\(\xi_l\), \(\xi_r\)},
					yticklabel=\empty,
					ytick=\empty,
					xmin=0,
					xmax=4,
					ymin=0,
					ymax=4,
					axis x line=middle,
					axis y line=middle,
					y axis line style={-stealth, thick},
					x axis line style={-stealth, thick}
				]

				\addplot[black, thick, mark=o] coordinates { (-0.1, 3) (1.5, 3) }
				node [above, midway] {\(u_l\)};

				\addplot[black, thick, mark=o, mark repeat={2}] coordinates { (1.5, 2) (2, 2) }
				node [above] {\(u_m\)};
				\addplot[black, thick, mark=*, mark repeat={2}, mark options={
							fill=black, draw=black,
						}] coordinates { (2.5, 2) (2, 2) };

				\addplot[black, thick, mark=*, mark options={
							fill=black, draw=black,
						}] coordinates { (2.5, 1) (4.1, 1) } node [above, midway] {\(u_r\)};
			\end{axis}
		\end{tikzpicture}
	\end{center}
	\caption{Initial field for the \(a=x\) example.}\label{fig:ch5}
\end{marginfigure}

The only difficulty lies in the more complex ODE, which we can now
write as
\[
	\odv{x_\star}{t} = x, \quad x_\star(0) = x_0.
\]
We know\footnote{Check!} that the solution to this ODE is
\[
	x_\star = x_0 e^t,
\]
so we again can draw the characteristics that go through the points with
discontinuities and \say{paint} the solution (see \cref{fig:ch6}).

\begin{marginfigure}
	\begin{center}
		\begin{tikzpicture}
			% \pgfplotsset{every axis plot post/.append style={draw=black, fill=gray!25}}
			\begin{axis}[
					smooth,
					% area style,
					xlabel=\(x\),
					ylabel={\(t\)},
					width=6.3cm,
					xtick={0.5, 1},
					xticklabels={\(\xi_l\), \(\xi_r\)},
					yticklabel=\empty,
					ytick=\empty,
					xmin=0,
					xmax=4,
					ymin=0,
					ymax=3,
					axis x line=middle,
					axis y line=middle,
					y axis line style={-stealth, thick},
					x axis line style={-stealth, thick}
				]
				\addplot[name path=D] coordinates {(1, -0.1) (4, -0.1)};
				\addplot[name path=U] coordinates {(0, 4.1) (4, 4.1)};

				\addplot[name path=A, black, thick, domain = 0.5:4,samples=20]  {ln(x/0.5)}
				node [midway, above, xshift=-0.8cm] {\(x_\star = \xi_l e^t\)} ;
				\addplot[name path=B, black, thick, domain = 1:4,samples=20]  {ln(x/1)}
				node [midway, below, xshift=0.8cm] {\(x_\star = \xi_r e^t\)} ;
				\fill [green!25] (0, 0) -- (0.5, 0) -- (0,4.1);
				\addplot[blue!25] fill between[of=A and B];
				\addplot[green!25] fill between[of=A and U];
				\addplot[red!25] fill between[of=B and D];
				\node at (2, 2.5) {\(u_l\)};
				\node at (3, 1.5) {\(u_m\)};
				\node at (3.75, 1) {\(u_r\)};
			\end{axis}
		\end{tikzpicture}
	\end{center}
	\caption{Solution to the \(a=x\) example.}\label{fig:ch6}
\end{marginfigure}

\section{Non-linear hyperbolic conservation laws}

\newthought{The Euler equations} are non-linear, due to the
tensor product\footnote{%
	Two different indices multiplied are a tensor product, i.e., if a tensor is
	formed from the product \(\bm{T} = \bm{u}\otimes\bm{v}\), then this can be
	written as \(\bm{T} = u_i u_j \bm{e}_i\otimes\bm{e}_j\) and \(\bm{e}_i\)
	denotes an \(i\)-th versor of the Cartesian coordinate system. From here, we
	can also deduce that \(T_{ij}= u_i u_j\).
}
of velocities under the divergence term. Consider
\[
	\pdv{\rho u_i}{t} + \pdv{\rho u_i u_j}{x_j} = -\pdv{p}{x_i}.
\]
Here the flux function is given as \(f(u_i) = \rho u_i u_j\).
The simplest model that would allow us to study such non-linear behaviour
is the Burgers equation, where \(f(u) = u^2/2\):
\begin{equation}
	\pdv{u}{t} + \pdv{u^2/2}{x} = 0,
	\label{eq:burgers}
\end{equation}
or in linearized form
\begin{equation}
	\pdv{u}{t} + u\pdv{u}{x} = 0.
	\label{eq:lin-burgers}
\end{equation}
To linearize, we computed the derivative \(\pdv{f}/{u}\) as in the previous
example.

\subsection{Emergence of discontinuities}
We know that the solution can be discontinuous if \(u_0\) is not continuous. In
case of non-linear equations, discontinuous solution can emerge even if the
initial condition is continuous. We will consider the following example, with
\(u_0\) given as a tent function centred around \(\xi\), as sketched in
\cref{fig:burg1}.

\begin{marginfigure}
	\begin{center}
		\begin{tikzpicture}
			% \pgfplotsset{every axis plot post/.append style={draw=black, fill=gray!25}}
			\begin{axis}[
					% area style,
					xlabel=\(x\),
					ylabel={\(u_0\)},
					width=6.3cm,
					height=3cm,
					xtick={1, 2, 3},
					xticklabels={\(\xi_l\), \(\xi\), \(\xi_r\)},
					yticklabel=\empty,
					ytick=\empty,
					xmin=0,
					xmax=4,
					ymin=0,
					ymax=3,
					axis x line=middle,
					axis y line=middle,
					y axis line style={-stealth, thick},
					x axis line style={-stealth, thick}
				]
				\addplot[black, thick] coordinates { (0, 0) (1, 0) (2, 2) (3, 0) (4,0) };
			\end{axis}
		\end{tikzpicture}
	\end{center}
	\caption{Initial field for the Burgers equation example.}\label{fig:burg1}
\end{marginfigure}

Before we try to draw characteristic curves there are a couple of points to
consider.
\begin{enumerate}
	\item The ODE is given as
	      \[
		      \odv{x_\star}{t} = u(x, t).
	      \]
	\item We know that \(\odv{u}/{t} = 0\) for the curve so, we can treat the
	      right-hand side of the above equation as constant.
	\item This means that for each point \(x_0\), the characteristic curve will
	      be a straight line \(x_\star = x_0 + u_0(x_0)t\), but the slope of the curves
	      can change depending on the value of \(u_0\).
\end{enumerate}


\begin{marginfigure}
	\begin{center}
		\begin{tikzpicture}
			% \pgfplotsset{every axis plot post/.append style={draw=black, fill=gray!25}}
			\begin{axis}[
					% area style,
					xlabel=\(x\),
					ylabel={\(t\), {\color{gray}\(u_0\)}},
					width=6.3cm,
					xtick={1, 2, 3},
					xticklabels={\(\xi_l\), \(\xi\), \(\xi_r\)},
					yticklabel=\empty,
					ytick=\empty,
					xmin=0,
					xmax=4,
					ymin=0,
					ymax=3,
					axis x line=middle,
					axis y line=middle,
					y axis line style={-stealth, thick},
					x axis line style={-stealth, thick}
				]
				\addplot[black, gray] coordinates { (0, 0) (1, 0) (2, 0.5) (3, 0) (4,0) };

				\addplot[red, thick] coordinates { (0.25, 0) (0.25, 2) };
				\addplot[red, thick] coordinates { (0.5, 0) (0.5, 2) };
				\addplot[red, thick] coordinates { (0.75, 0) (0.75, 2) };
				\addplot[red, thick] coordinates { (1, 0) (1, 2) };

				\addplot[red, thick, domain=1:1.5] {(x - 1.25)/0.125};
				\addplot[red, thick, domain=1:2] {(x - 1.5)/0.25};
				\addplot[red, thick, domain=1:2.5] {(x - 1.75)/0.375};

				\addplot[red, thick, domain=2:3] {(x - 2)/0.5};
				\addplot[red, thick, domain=2:3] {(x - 2.25)/0.375};
				\addplot[red, thick, domain=2:3] {(x - 2.5)/0.25};
				\addplot[red, thick, domain=2:3] {(x - 2.75)/0.125};

				\addplot[red, thick] coordinates { (3, 0)    (3, 2) };
				\addplot[red, thick] coordinates { (3.25, 0) (3.25, 2) };
				\addplot[red, thick] coordinates { (3.5, 0)  (3.5, 2) };
				\addplot[red, thick] coordinates { (3.75, 0) (3.75, 2) };

				\fill    [green!25] (0, 2) -- (4,2) -- (4, 3) -- (0, 3);
				\addplot [dashed, thick] coordinates { (0, 2) (4, 2) } node [above, midway] {\(t_d\) line, discontinuity above};


				% \addplot[red, thick] {-2.5 + 0.25*x};
			\end{axis}
		\end{tikzpicture}
	\end{center}
	\caption{Solution to the Burgers equation until a time a discontinuity appears.}\label{fig:burg2}
\end{marginfigure}

\newthought{The solution} can be obtained by dealing with the
initial condition in intervals:
\begin{description}
	\item[\(x_0 < \xi_l\):] here \(u_0 = 0\) which means that characteristic is
	      given by \(x_\star = x_0\) (i.e., straight vertical line going through \(x_0\));
	\item[\(x_0 \geq \xi_r\):] same as above;
	\item[\(\xi_l \geq x_0 < \xi\):] as the value of \(u_0\) increases in this interval, from
	      0 to a maximum value, causing the characteristics to flatten progressively and creating the
	      \say{fan} pattern;
	\item[\(\xi \geq x_0 < \xi+r\):] as the value of \(u_0\) decreases in this interval, from
	      maximum value to 0, causing the characteristics to steepen and cross in one point!
\end{description}
We can plot it (see \cref{fig:burg2}), and we see that at time \(t=t_d\), sever
curves cross. This means, that this point has more than one value (as the value
on the characteristic curve remains constant, and we have conflicting values on
different curves). This means, that the solution \(t>t_d\) is \emph{discontinuous},
even though \(u_0 \in C^0\). We can't say anything about the solution for \(t>t_d\)
using the method of characteristics.

\section{Weak solutions for non-linear hyperbolic PDEs}
What does it mean for the problem \cref{eq:burgers}, that the solution at a
certain time is discontinuous? Importantly, this means that \emph{the
	derivatives do not exist}. A solution has to be defined in a different
sense, without requiring that \(\pdv{u}/{x}\) and \(\pdv{u}/{t}\) exist.

We can define a test function \(\Phi(x,t)\in
C^1(\mathbb{R}\times\mathbb{R}^+)\), which is differentiable, and its value
vanishes at infinite (\(\Phi(-\infty, t) = 0\), \(\Phi(\infty, t) = 0\),
\(\Phi(x, \infty) = 0\)).

The weak formulation of \cref{eq:burgers} is derived by multiplying the equation by
\(\Phi\)
\[
	\Phi\pdv{u}{t} + \Phi\pdv{f}{x} = 0,
\]
and integrating in time and space
\[
	\int\limits_{-\infty}^{\infty} \int\limits_{0}^{\infty} \Phi\pdv{u}{t} \d t \d x
	+ \int\limits_{-\infty}^{\infty} \int\limits_{0}^{\infty} \Phi\pdv{f}{x}  \d t \d x= 0.
\]
Now we can use integration by parts\footnote{%
	A quick reminder, for functions \(u(x)\) and \(v(x)\), the definite integral on the
	interval \([a,b]\) can be reformulated as
	\[
		\int\limits_a^b \odv{u}{x} v \d x =
		uv|_a^b
		- \int\limits_a^b u\odv{v}{x}  \d x ,
	\]
	where
	\[
		uv|_a^b = u(b)v(b) - u(a)v(a).
	\]
}
on the parts of the above integrals. Dealing first with the time derivative term
\begin{equation*}
	\begin{split}
		\int\limits_{-\infty}^{\infty} \int\limits_{0}^{\infty} \Phi\pdv{u}{t} \d t \d x
		= & -\int\limits_{-\infty}^{\infty} \int\limits_{0}^{\infty} \pdv{\Phi}{t}u \d t \d x   \\
		  & + \int\limits_{-\infty}^{\infty} \underbrace{u(x, \infty)\Phi(x, \infty)}_{=0} \d x
		- \int\limits_{-\infty}^{\infty} u(x, 0)\Phi(x, 0) \d x.
	\end{split}
\end{equation*}
In case of the second term we can reverse the order of integration and do the same thing:
\begin{equation*}
	\begin{split}
		\int\limits_{0}^{\infty} \int\limits_{-\infty}^{\infty} \Phi\pdv{f}{x}  \d x \d t
		= & -\int\limits_{0}^{\infty} \int\limits_{-\infty}^{\infty} \pdv{\Phi}{x}f \d x \d t \\
		  & + \int\limits_{0}^{\infty} \underbrace{f(\infty, t)\Phi(\infty, t)}_{=0} \d t
		- \int\limits_{0}^{\infty} \underbrace{f(-\infty, t)\Phi(-\infty, t)}_{=0} \d t.
	\end{split}
\end{equation*}
If we resemble the full expression we get the weak formulation of \cref{eq:burgers}
\[
	\int\limits_{-\infty}^{\infty} \int\limits_{0}^{\infty} \pdv{\Phi}{t}u + \pdv{\Phi}{x}f \d t \d x
	+ \int\limits_{-\infty}^{\infty} u(x, 0)\Phi(x, 0) \d x
	= 0
\]
The above equation is equivalent because \(\Phi\) can be chosen arbitrarily,\footnote{%
	That's a quite complicated topic, look into the course literature if you are
	interested.}
but importantly, there are no derivatives of \(u\) in the weak form. Hence, we can
use it to try to reason about the solution with discontinuities.

\subsection{Rankine--Hugoniot formula (shock front speed)}
We return to a previous problem with initial condition given
as
\[
	u_0 = \left\{
	\begin{aligned}
		u_l \; & : \; x < \xi     \\
		u_r \; & : \; x \geq \xi,
	\end{aligned} \right.
\]
and possibly non-linear equation
\[
	\pdv{u}{t} + \pdv{f(u)}{x} = 0.
\]
The discontinuity is already present in the solution, so we are looking
for the solution in the weak sense.

\begin{marginfigure}
	\begin{center}
		\begin{tikzpicture}
			% \pgfplotsset{every axis plot post/.append style={draw=black, fill=gray!25}}
			\begin{axis}[
					smooth,
					% area style,
					xlabel=\(x\),
					ylabel=\empty,
					width=6.3cm,
					height=3cm,
					xtick={0.5, 1.5},
					xticklabels={\(\xi\), \(\xi_1\)},
					yticklabel=\empty,
					ytick=\empty,
					xmin=0,
					xmax=4,
					ymin=0,
					ymax=4,
					axis x line=middle,
					axis y line=middle,
					y axis line style={-stealth, thick},
					x axis line style={-stealth, thick}
				]

				\addplot[dashed, thick, red, shift={(1,0)}] coordinates { (0.5, 2.5) (0.5, 0.5)};
				\addplot[red, thick, shift={(1, 0)}] coordinates { (-2, 2.5) (0.5, 2.5) }  node [above] {\(u(x, t_1)\)};
				\addplot[red, thick, shift={(1, 0)}] coordinates { (0.5, 0.5) (4.1, 0.5) } ;

				\draw [-stealth, thick, blue] (1.5, 2) -- (2.5, 2) node [midway, above] {\(s\)};

				\addplot[dashed, thick, black] coordinates { (0.5, 2.5) (0.5, 0.5)};
				\addplot[black, thick, mark=o] coordinates { (-0.1, 2.5) (0.5, 2.5) }
				node [below, midway] {\(u_l\)};
				\addplot[black, thick, mark=*, mark options={
							fill=black, draw=black,
						}] coordinates { (0.5, 0.5) (4.1, 0.5) } node [above, midway] {\(u_r\)};
			\end{axis}
		\end{tikzpicture}
	\end{center}
	\caption{Initial field for the Rankine--Hugoniot problem example along with the solution at time \(t_1\).}\label{fig:rk1}
\end{marginfigure}

The solution is of the form (see \cref{fig:rk1})
\[
	u(x,t) = \left\{
	\begin{aligned}
		u_l \; & : \; x < \xi + st,    \\
		u_r \; & : \; x \geq \xi + st,
	\end{aligned} \right.
\]
where \(s\) is the speed of the shock front, that we have to determine. Let \(M\) be a large number.
We can integrate the PDE
\[
	\int\limits_{-M}^{M} \pdv{u}{t} \d x = -  \int\limits_{-M}^{M} \pdv{f}{x} \d x
\]
and compute the right-hand side directly
\[
	- \int\limits_{-M}^{M} \pdv{f}{x} \d x  = -f(u(M, t)) + f(u(-M, t)) = f(u_l) - f(u_r).
\]
Similarly, for the time derivative
\[
	\int\limits_{-M}^{M} \pdv{u}{t} \d x
	= \pdv{}{t}\int\limits_{-M}^{M} u \d x
	= \pdv{}{t}\ab\{ (M+st)u_l + (M-st)u_r \} = s(u_l - u_r)
\]
We can equate both expressions
\[
	s(u_l - u_r) = f(u_l) - f(u_r)
\]
to get the \emph{Rankine--Hugoniot} formula for the speed of the shock front:
\[
	s = \frac{f(u_l) - f(u_r)}{u_l - u_r}.
\]

% \begin{fullwidth}
% \bibliography{refs}
% \end{fullwidth}


\end{document}
